\chapter{Testing}
\label{ch:testing}

The correctness of TalkCS was verified with automated tests on both the backend and the
frontend, complemented by manual testing of the complete application in the browser.
This chapter describes the testing strategy, the test suites, and their results.

\section{Testing Strategy}

The automated tests follow the idea of the \emph{test pyramid}: a large number of fast,
isolated unit tests for the business logic, a smaller number of tests for the web and
persistence layers, and a few end-to-end style integration tests that exercise the
complete stack. Table~\ref{tab:test-levels} gives an overview.

\begin{table}[H]
  \centering
  \small
  \begin{tabularx}{\textwidth}{@{}l X r r@{}}
    \toprule
    \textbf{Level} & \textbf{Tools and scope} & \textbf{Classes / files} & \textbf{Tests} \\
    \midrule
    Backend service unit tests & JUnit 5, Mockito; services with mocked repositories & 20 & 118 \\
    Backend controller tests & \code{@WebMvcTest}, MockMvc; HTTP layer with mocked services & 14 & 57 \\
    Backend repository tests & \code{@DataJpaTest}, H2; JPA queries & 8 & 25 \\
    Backend integration tests & \code{@SpringBootTest}; full application context & 3 & 5 \\
    \midrule
    Frontend API client tests & Vitest; Axios mocked & 14 & 60 \\
    Frontend page, hook and component tests & Vitest, React Testing Library, jsdom & 13 & 40 \\
    \bottomrule
  \end{tabularx}
  \caption{Overview of the automated test suites}
  \label{tab:test-levels}
\end{table}

All tests can be run with a single command from the project root:

\begin{lstlisting}[language=bash,caption={Running all tests},label={lst:testsh}]
./test.sh
# which is equivalent to:
cd backend  && ./gradlew test          # backend tests + JaCoCo report
cd frontend && npm test -- --run       # frontend tests
\end{lstlisting}

Backend tests use a separate configuration profile (\code{application-test.properties})
which replaces PostgreSQL with an H2 in-memory database in PostgreSQL compatibility mode,
uses a fixed test JWT secret, and a dummy e-mail API key, so the tests never touch the
development database or send real e-mails. Each test that needs a logged-in user sets the
security context through a small helper class (\code{SecurityTestSupport}).

\section{Backend Unit Tests}

Service tests are plain JUnit 5 tests that run with the \code{MockitoExtension}. The
service under test is created with \code{@InjectMocks}, and all of its dependencies
(repositories, other services, the messaging template) are replaced with Mockito mocks.
These tests run in milliseconds and check the business rules in isolation, both the
expected behaviour and the error cases. The most important tested rules are:

\begin{description}
  \item[AuthService] registration hashes the password and returns a token; duplicate
    e-mail addresses and usernames are rejected; login fails for unknown users and wrong
    passwords.
  \item[PasswordResetService] unknown e-mail addresses are silently ignored; a valid
    request stores a token and sends an e-mail whose link is built from the configured
    base URL; invalid, used and expired tokens are
    rejected; a successful reset updates the password and marks the token as used.
  \item[VoteService] self-votes are blocked; a new vote changes the author's reputation;
    clicking the same vote again removes it; changing the direction adjusts the
    reputation by the difference.
  \item[PostService] (13 tests) posts are created and linked to existing categories only;
    only the author can edit; the author and administrators can delete; accepting an
    answer updates the reputation and notifies the answerer, and is blocked for other
    users; withdrawing an acceptance reverses the rewards; bookmarking is idempotent.
  \item[BadgeService] the first-post badge is awarded at the threshold and not below it;
    badges are never awarded twice; the bronze expertise badge is awarded at 50 category
    reputation.
  \item[SimilarityService] an empty category yields no suggestions; a post sharing
    specific terms with the draft is ranked above unrelated posts.
  \item[TagSuggestionService] matching tags are returned; stop words and short tokens are
    ignored; at most five tags are suggested.
  \item[MentionService] mentions are extracted, de-duplicated and restricted to existing
    users.
  \item[ChatService] conversation participants are stored in normalized order and
    reused; messages are persisted and broadcast to both participants; non-members cannot
    send messages.
  \item[CalendarService] events of administrators are public and those of students
    private; proposals are created as pending; approval creates a public event; rejection
    stores the note; non-owners cannot edit events.
  \item[EmailDigestService] digests are sent only to users with recent unread
    notifications who have not opted out.
\end{description}

Code~\ref{lst:test-vote} shows a typical service test, which checks that withdrawing a
vote restores the author's reputation.

\begin{lstlisting}[language=Java,caption={Unit test of vote toggling (\texttt{VoteServiceTests.java})},label={lst:test-vote}]
@Test
void voteOnPost_togglingSameVoteRemovesIt() {
    SecurityTestSupport.setAuth("voter@t.com");
    User voter = user(2L, "voter@t.com", 0);
    User author = user(1L, "a@t.com", 11);
    Post p = Post.builder().id(1L).author(author).category(cat()).build();
    Vote existing = Vote.builder().id(50L).voter(voter).post(p).value(1).build();
    when(userRepository.findByEmail("voter@t.com")).thenReturn(Optional.of(voter));
    when(postRepository.findById(1L)).thenReturn(Optional.of(p));
    when(voteRepository.findByVoterIdAndPostId(2L, 1L)).thenReturn(Optional.of(existing));
    when(categoryReputationRepository.findByUserIdAndCategoryId(any(), any()))
        .thenReturn(Optional.empty());

    voteService.voteOnPost(1L, 1);
    assertThat(author.getReputation()).isEqualTo(10);
    verify(voteRepository).delete(existing);
}
\end{lstlisting}

\section{Backend Controller Tests}

Controller tests use Spring's \code{@WebMvcTest} slice, which starts only the web layer
of the application: the controller under test, JSON serialization, validation and
exception handling. Services are replaced with \code{@MockitoBean} mocks, and requests are
sent with \code{MockMvc} without starting a real server. Every one of the fourteen REST
controllers has its own test class, which verifies the URL mappings, request parameters,
status codes and the structure of the JSON responses.

Three of the classes (\code{AdminControllerTests}, \code{CategoryControllerTests} and
\code{CalendarControllerTests}) additionally import a configuration that enables method
security, and use \code{@WithMockUser} to verify the role-based access rules. For example,
they check that a user with the Student role receives \code{403 Forbidden} when trying to
read the admin statistics, create or archive a category, list the pending calendar
proposals or approve a proposal (Code~\ref{lst:test-403}).

\begin{lstlisting}[language=Java,caption={Verifying role-based access control (\texttt{CategoryControllerTests.java})},label={lst:test-403}]
@Test
@WithMockUser(roles = "STUDENT")
void createCategory_returns403ForStudent() throws Exception {
    CategoryRequest req = new CategoryRequest();
    req.setName("X"); req.setDescription("d");
    mockMvc.perform(post("/api/categories")
            .contentType(MediaType.APPLICATION_JSON)
            .content(objectMapper.writeValueAsString(req)))
        .andExpect(status().isForbidden());
}
\end{lstlisting}

\section{Backend Repository Tests}

Repository tests use the \code{@DataJpaTest} slice, which starts Hibernate with the H2
database and rolls back every test in a transaction. They verify the derived and custom
queries that the services rely on, for example that archived categories are excluded,
that only root comments are returned for a post, that posts sorted by vote score or
comment count are ordered correctly across pages, that post and comment bodies longer
than 255 characters are stored, that unread notifications are counted
and marked as read correctly, that votes are filtered by value, and that the popular
tags query orders tags by the number of posts using them.

\section{Backend Integration Tests}

Three test classes start the complete Spring application context with
\code{@SpringBootTest}. The first one simply checks that the application context loads, which catches configuration
errors such as missing beans or invalid properties. The second one,
\code{UserLifecycleIntegrationTests}, exercises the whole security chain through real HTTP
requests: it registers a new user, logs in with the same credentials, extracts the JWT
from the response, and calls the protected \code{GET /api/users/me} endpoint with the
token, expecting a successful response. This test passes through the real
\code{JwtAuthFilter}, \code{UserLoader}, services and repositories. The third one,
\code{CorsIntegrationTests}, sends CORS preflight requests and checks that the configured
frontend URL and the local development server are allowed, while any other origin is
rejected.

\section{Backend Test Results and Coverage}

All 205 backend tests in 45 test classes pass. The complete suite runs in about ten
seconds. Code coverage is measured with JaCoCo, which is configured in Gradle to produce
an HTML report after every test run (\code{backend/build/reports/jacoco/test/html}).
Table~\ref{tab:coverage} summarizes the coverage per package.

\begin{table}[H]
  \centering
  \small
  \begin{tabular}{@{}l r r@{}}
    \toprule
    \textbf{Package} & \textbf{Instruction coverage} & \textbf{Branch coverage} \\
    \midrule
    \code{config} & 95\% & 50\% \\
    \code{model} & 90\% & 0\% \\
    \code{dto} & 78\% & 50\% \\
    \code{service} & 70\% & 63\% \\
    \code{controller} & 52\% & 0\% \\
    \code{security} & 45\% & 13\% \\
    \midrule
    \textbf{Total} & \textbf{69\%} & \textbf{57\%} \\
    \bottomrule
  \end{tabular}
  \caption{Backend code coverage measured with JaCoCo}
  \label{tab:coverage}
\end{table}

The business logic in the \code{service} package, which contains most of the decisions
of the application, is the most thoroughly tested part. The lower coverage of the
\code{security} package is explained by the fact that the controller tests deliberately
disable the security filters to test the web layer in isolation; the authentication
filter is exercised by the integration test instead.

\screenshot{backend-tests}{Gradle test report of the backend test suite}{fig:backend-tests}{1.0}

\section{Frontend Tests}

Frontend tests are written with Vitest, which shares its configuration with Vite, and run
in a simulated browser environment (jsdom).

\subsubsection*{API Client Tests}

Each module of the \code{api} directory has a test file in which the shared Axios
instance is replaced by a mock (\code{vi.mock('./api')}). The tests verify that every
function calls the correct HTTP method and URL, encodes the query parameters correctly
(for example, that \code{getPosts} omits \code{categoryId} when it is not given), sends
the expected request body, and returns the data of the response. These tests guard the
contract between the frontend and the REST API described in Section~\ref{sec:architecture}.

\subsubsection*{Page, Hook and Component Tests}

Components are rendered with React Testing Library and tested from the user's
perspective, by finding elements by their label or text and simulating user input.
Table~\ref{tab:fe-tests} lists the covered components.

\begin{table}[H]
  \centering
  \small
  \begin{tabularx}{\textwidth}{@{}l X@{}}
    \toprule
    \textbf{Component} & \textbf{Tests covered} \\
    \midrule
    \code{LoginPage} & Renders inputs; validation errors on empty submit; calls the login API on valid input; shows an error on failed login \\
    \code{RegisterPage} & Renders fields; minimum length error for short passwords; calls the register API on valid input \\
    \code{ForgotPasswordPage} & Renders input; success message after submit; error on failure \\
    \code{ResetPasswordPage} & Invalid link without token; rejects mismatched and short passwords; submits valid token and password \\
    \code{SearchPage} & Prompt for empty query; calls search and shows the result count; no-results state \\
    \code{BadgesPage} & Fetches and renders badges; empty badge list \\
    \code{LeaderboardPage} & Loading state; renders heading and users \\
    \code{ProtectedRoute} & Redirects unauthenticated users to \route{/login}; renders children with a token \\
    \code{useAuth} & Default state; login stores the token; logout clears the storage \\
    \code{usePostUpdates} & No connection without category or token; subscribes on connect; dispatches created, updated and deleted events; deactivates on unmount \\
    \code{useCommentUpdates} & No connection without post or token; refreshes when a frame arrives; deactivates on unmount \\
    \code{Button}, \code{utils} & Variants and class name merging \\
    \bottomrule
  \end{tabularx}
  \caption{Frontend page, hook and component test coverage}
  \label{tab:fe-tests}
\end{table}

The WebSocket hooks are tested by mocking the STOMP client: the tests capture the
\code{onConnect} callback, invoke it manually, deliver fake message frames, and assert
that the hook subscribes to the right destination, forwards the events, and deactivates
the client when the component unmounts. In total the frontend suite contains 100 tests in
27 files, and all of them pass.

\screenshot[!htb]{frontend-tests}{Vitest output of the frontend test suite}{fig:frontend-tests}{0.85}

\section{Manual Testing}

After each development phase, the complete application was started locally and the main
user flows were tested manually: the backend endpoints with \code{curl}, and the user
interface in the browser. The manual tests focused on aspects that are hard to automate:

\begin{itemize}
  \item real-time behaviour with two browsers logged in as different users (live comments,
    live post lists, chat delivery and unread counters);
  \item the role-dependent user interface for Student, Professor and Admin accounts;
  \item the responsive layout on narrow screens using the browser's device emulation;
  \item the dark and light themes on every page;
  \item file uploads and downloads of different types and sizes, including the 50~MB
    limit;
  \item the password reset e-mail and link, and the delivery of notification e-mails.
\end{itemize}

Errors found during manual testing were fixed before the corresponding change was
committed to the repository.
